\documentclass[11pt]{article}

\usepackage[margin=1in]{geometry}
\usepackage{setspace}
\usepackage{graphicx}
\usepackage{booktabs}
\usepackage[colorlinks=true, linkcolor=blue, urlcolor=blue]{hyperref}
\usepackage{float}
\usepackage{xcolor}
\usepackage{amsmath}

\graphicspath{{output/graphs/descriptives/}{output/tables/descriptives/}}

\title{Descriptive Statistics\\
       \large AI Exposure \& Labor Market Outcomes --- March 30, 2026}
\author{Leyao Wang}
\date{}

\begin{document}
\maketitle
\thispagestyle{empty}

\noindent\textit{Sample:} CPS-ASEC 2018--2025, labor-force participants aged 22--65 with at least a high school diploma ($N\approx 500{,}669$). AI exposure is the Eloundou et al.\ (2024) GPT-4 $\beta$ score (\texttt{dv\_rating\_beta}), ranging from 0 to $\approx$0.78, assigned at the occupation level and time-invariant. \textit{College} = BA or above (educ $>$ 110). \textit{Non-college} = high school or some college (educ $\leq$ 110).

\bigskip

%==================================================
\section*{Contents}
%==================================================

\subsection*{Overall distribution}
\begin{enumerate}
  \item \hyperref[fig:kdensity_overall]{Kernel density of AI exposure across workers}
  \item \hyperref[tab:top5]{Most exposed occupations (top 5)}
  \item \hyperref[tab:bot5]{Least exposed occupations (bottom 5)}
\end{enumerate}

\subsection*{Heterogeneity by education}
\begin{enumerate}\setcounter{enumi}{3}
  \item \hyperref[fig:kdensity_educ]{Kernel density of AI exposure: college vs.\ non-college}
\end{enumerate}

\subsection*{Time-series and summary statistics}
\begin{enumerate}\setcounter{enumi}{4}
  \item \hyperref[fig:unemp_quartile]{Unemployment rate by AI exposure quartile (post-2022)}
  \item \hyperref[tab:summary]{Summary statistics by AI exposure quartile}
\end{enumerate}

\clearpage
\doublespacing

%==================================================
\section{Overall Distribution of AI Exposure}
%==================================================

\begin{figure}[H]
\begin{center}
\includegraphics[width=0.88\textwidth]{kdensity_beta_overall.pdf}
\caption{Kernel density of AI exposure ($\beta_o$) across workers in the CPS-ASEC sample. Each observation is a worker-year; the density reflects the worker-level distribution of occupational AI exposure scores.}
\label{fig:kdensity_overall}
\end{center}
\smallskip
\noindent\footnotesize\textit{Notes:} AI exposure is the Eloundou et al.\ (2024) GPT-4 $\beta$ score, assigned at the occupation level. Higher values indicate greater exposure to LLM automation. The score is time-invariant; variation is entirely cross-occupational.
\end{figure}

\begin{table}[H]
\begin{center}
\caption{Most exposed occupations (top 5 by $\beta_o$).}
\label{tab:top5}
\begin{tabular}{lc}
\hline
Occupation & Exposure Score \\
\hline
Mathematicians&1.000\\
Proofreaders and Copy Markers&0.975\\
Blockchain Engineers&0.971\\
Correspondence Clerks&0.964\\
Court Reporters and Simultaneous Captioners&0.958\\
\hline
\end{tabular}

\end{center}
\smallskip
\noindent\footnotesize\textit{Source:} Eloundou et al.\ (2024), occupation-level file (\texttt{occ\_level.csv}). Scores are the GPT-4-annotated $\beta$ measure.
\end{table}

\begin{table}[H]
\begin{center}
\caption{Least exposed occupations (bottom 5 by $\beta_o$).}
\label{tab:bot5}
\begin{tabular}{lc}
\hline
Occupation & Exposure Score \\
\hline
Hoist and Winch Operators&0.000\\
Continuous Mining Machine Operators&0.000\\
Derrick Operators, Oil and Gas&0.000\\
Packaging and Filling Machine Operators and Tenders&0.000\\
Rail-Track Laying and Maintenance Equipment Operators&0.000\\
\hline
\end{tabular}

\end{center}
\smallskip
\noindent\footnotesize\textit{Source:} Eloundou et al.\ (2024), occupation-level file (\texttt{occ\_level.csv}). Scores are the GPT-4-annotated $\beta$ measure.
\end{table}

\clearpage

%==================================================
\section{AI Exposure by Education}
%==================================================

\begin{figure}[H]
\begin{center}
\includegraphics[width=0.88\textwidth]{kdensity_beta_education.pdf}
\caption{Kernel density of AI exposure ($\beta_o$): college-educated workers (BA or above, navy solid) vs.\ non-college workers (high school or some college, maroon dashed).}
\label{fig:kdensity_educ}
\end{center}
\smallskip
\noindent\footnotesize\textit{Notes:} Each observation is a worker-year. The two distributions show the occupational exposure profile faced by each education group. If the distributions overlap substantially, the asymmetric wage result (non-college compressed, college not) is more striking because it cannot be attributed simply to differential exposure.
\end{figure}

\clearpage

%==================================================
\section{Time-Series and Summary Statistics}
%==================================================

\begin{figure}[H]
\begin{center}
\includegraphics[width=0.88\textwidth]{unemp_by_quartile_post2022.pdf}
\caption{Mean unemployment rate by AI exposure quartile, 2022--2025. Workers are divided into four groups based on quartiles of the AI exposure score ($\beta_o$). Q1 (lowest exposure) to Q4 (highest exposure).}
\label{fig:unemp_quartile}
\end{center}
\smallskip
\noindent\footnotesize\textit{Notes:} Raw means, no regression adjustment. The figure shows whether high-exposure workers experienced differential unemployment trends in the post-ChatGPT period. A divergence between Q4 and Q1 after 2022 is the ``eyeball'' version of the formal DiD estimates.
\end{figure}

\begin{table}[H]
\begin{center}
\caption{Summary statistics by AI exposure quartile.}
\label{tab:summary}
\begin{table}[htbp]\centering
\def\sym#1{\ifmmode^{#1}\else\(^{#1}\)\fi}
\caption{Summary Statistics by AI Exposure Quartile}
\begin{tabular}{l*{5}{cc}}
\toprule
                    &\multicolumn{2}{c}{All}  &\multicolumn{2}{c}{Q1 (low)}&\multicolumn{2}{c}{Q2}   &\multicolumn{2}{c}{Q3}   &\multicolumn{2}{c}{Q4 (high)}\\
                    &        mean&          sd&        mean&          sd&        mean&          sd&        mean&          sd&        mean&          sd\\
\midrule
Age                 &      42.312&      11.749&      41.110&      12.226&      42.060&      11.822&      43.584&      11.087&      42.544&      11.672\\
College (BA+)       &       0.442&       0.497&       0.153&       0.360&       0.479&       0.500&       0.598&       0.490&       0.553&       0.497\\
Log Real Wage       &       9.733&       0.928&       9.347&       0.885&       9.583&       0.873&      10.086&       0.874&       9.917&       0.896\\
Unemployed          &       0.034&       0.181&       0.055&       0.228&       0.031&       0.173&       0.022&       0.145&       0.026&       0.160\\
AI Exposure ($\beta\_o$)&       0.370&       0.208&       0.096&       0.063&       0.318&       0.047&       0.452&       0.029&       0.627&       0.122\\
\midrule
Observations        &      500669&            &      129090&            &      123745&            &      123871&            &      123963&            \\
\bottomrule
\end{tabular}
\end{table}

\end{center}
\smallskip
\noindent\footnotesize\textit{Notes:} Sample: CPS-ASEC 2018--2025, labor-force participants aged 22--65 with at least a high school diploma. Quartiles defined by the Eloundou et al.\ (2024) GPT-4 $\beta$ score. ``All'' is the pooled sample. Mean and standard deviation reported for each variable.
\end{table}

\end{document}
